\documentclass[A4]{twentysecondcv} % a4paper for A4 and letterpaper

%----------------------------------------------------------------------------------------
%	 PERSONAL INFORMATION
%----------------------------------------------------------------------------------------

% If you don't need one or more of the below, just remove the content leaving the command, e.g. \cvnumberphone{}

\profilepic{profile.png} % Profile picture

\cvname{Akansha Kumari} % Your name
\cvjobtitle{RD Professional} % Job title/career

\cvdate{} % Date of birth
\cvaddress{Dumka, Jharkhand} % Short address/location, use \newline if more than 1 line is required
\cvnumberphone{+91 6206163544} % Phone number
\cvsite{} % Personal website
\cvmail{Kr.akansha.rd@gmail.com} % Email address

%----------------------------------------------------------------------------------------

\begin{document}

%----------------------------------------------------------------------------------------
%	 ABOUT ME
%----------------------------------------------------------------------------------------
\aboutme{Social development professional with experience in livelihood promotion and women and child health programs. Skilled in beneficiary identification, livelihood planning, community mobilization, and implementation of government-supported schemes to enhance income and social welfare. Committed to advancing sustainable development in agriculture, health, sanitation, education, and rural livelihoods within community-based organizations.}
 % To have no About Me section, just remove all the text and leave \aboutme{}

%----------------------------------------------------------------------------------------
%	 SKILLS
%----------------------------------------------------------------------------------------

% Skill bar section, each skill must have a value between 0 an 6 (float)
\skills{{Livelihood\;Planning/5},{Community\;Mobilization/5.5},{Leadership/5.8}}

%------------------------------------------------

% Skill text section, each skill must have a value between 0 an 6
% \skillstext{{lovely/4},{narcissistic/3}}


\languages{{Nagpuri/4},{Bengali/4},{Hindi/6},{English/5}}


%----------------------------------------------------------------------------------------

\makeprofile % Print the sidebar







%----------------------------------------------------------------------------------------
%	 INTERESTS
%----------------------------------------------------------------------------------------

\section{Interests}

Community-led development, women’s empowerment, sustainable rural livelihoods, public health awareness, and resilient agricultural practices.
%----------------------------------------------------------------------------------------
%	 EDUCATION
%----------------------------------------------------------------------------------------

\section{Education}

\begin{twenty}

	\twentyitem
	{2022}
	{Integrated M.Sc. Agriculture, Rural and Tribal Development, RKMVERI, Ranchi (Marks: 80.8\%)}
	{}
	{\emph{Thesis:} “Effect of different organics on powdery mildew of pea (\textit{Erysiphe pisi}).”
	 \emph{Project:} “Natural Resource Use Efficiency of Small Holding Tribal Farmers in Ranchi District, Jharkhand.”}

	\twentyitem
	{2017}
	{I.Sc., CBSE, Ranchi (Marks: 62.4\%)}
	{}
	{Physics, Chemistry, Mathematics.}

    \twentyitem
    {2015}
    {Matriculation, JAC, Ranchi (Marks: 80.6\%)}
    {}
    {IT, Science, Social Science, English, Hindi.}

\end{twenty}

%----------------------------------------------------------------------------------------
%	 PUBLICATIONS
%----------------------------------------------------------------------------------------

\section{Publications/Seminars}

\begin{twenty} % Environment for a short list with no descriptions
    \twentyitem{2026}{AI in rural governance and Welfare}{}{Presented a paper on at the International Seminar on “Responsible AI: Towards Ethics, Governance, Scientific and Global Framework” organized by S.R.T College, Dhamri, Godda}
    \twentyitem
    {2022}
    {Kumari Akansha and Dr. Mishra Pankaj Kumar}
    {}
    {Effect of different organics on powdery mildew of pea (\textit{Erysiphe pisi}). 
    \textit{The Pharma Innovation Journal}, 2022; 11(11): 1760--1762. 
    Available at: \url{https://www.thepharmajournal.com/archives/2022/vol11issue11/PartV/11-11-89-613.pdf}.}
	%\twentyitem{<dates>}{<title/description>}
\end{twenty}

%----------------------------------------------------------------------------------------
%	 AWARDS
%----------------------------------------------------------------------------------------

\section{Awards}

\begin{twenty}
    \twentyitem
	{2021}
	{Youth for Change Award}
	{}
	{Awarded by Jai Jharkhand Abhiyan for work during rural living.}

	\twentyitem
	{2020}
	{Best Intern Award}
	{}
	{Awarded by Buland Udaan Organization for best work on menstrual hygiene.}
\end{twenty}

%----------------------------------------------------------------------------------------
%	 EXPERIENCE
%----------------------------------------------------------------------------------------

\section{Experience}

\begin{twenty}

	\twentyitem
	{Dec 2023 -- Oct 2026}
	{Block Lead, The/Nudge Institute (Under JSLPS)}
	{}
	{Worked under Phulo Jhano Ashirvad Abhiyan Project, supporting livelihood development through beneficiary identification, livelihood planning, and loan linkage via CLFs. Trained and supervised NJS cadres, monitored progress, verified data, facilitated access to welfare schemes, and promoted awareness on education, health, nutrition, and income enhancement.}

	\twentyitem
	{Jul 2022 -- Dec 2023}
	{Women and Child Health Coordinator, Karuna Shechen, India, Hata, East-Singhbhum}
	{}
	{Conducted community awareness on water-borne diseases, menstrual health, and maternal nutrition; mobilized and trained field workers; prepared IEC materials; managed data entry and analysis; supervised vermi-compost production at the organization’s kitchen garden.}

	\twentyitem
	{Dec 2020 -- Dec 2021}
	{Event Manager cum President (I/c), Gramin Bharti Seva Sanstha, Ranchi}
	{}
	{Organized and coordinated institutional events on a voluntary basis.}

\end{twenty}
%----------------------------------------------------------------------------------------
%	 OTHER INFORMATION
%----------------------------------------------------------------------------------------



%----------------------------------------------------------------------------------------
%	 SECOND PAGE EXAMPLE
%----------------------------------------------------------------------------------------

%\newpage % Start a new page

%\makeprofile % Print the sidebar

%\section{Other information}

%\subsection{Review}

%Alice approaches Wonderland as an anthropologist, but maintains a strong sense of noblesse oblige that comes with her class status. She has confidence in her social position, education, and the Victorian virtue of good manners. Alice has a feeling of entitlement, particularly when comparing herself to Mabel, whom she declares has a ``poky little house," and no toys. Additionally, she flaunts her limited information base with anyone who will listen and becomes increasingly obsessed with the importance of good manners as she deals with the rude creatures of Wonderland. Alice maintains a superior attitude and behaves with solicitous indulgence toward those she believes are less privileged.

%\section{Other information}

%\subsection{Review}

%Alice approaches Wonderland as an anthropologist, but maintains a strong sense of noblesse oblige that comes with her class status. She has confidence in her social position, education, and the Victorian virtue of good manners. Alice has a feeling of entitlement, particularly when comparing herself to Mabel, whom she declares has a ``poky little house," and no toys. Additionally, she flaunts her limited information base with anyone who will listen and becomes increasingly obsessed with the importance of good manners as she deals with the rude creatures of Wonderland. Alice maintains a superior attitude and behaves with solicitous indulgence toward those she believes are less privileged.

%----------------------------------------------------------------------------------------

\end{document}